% Modernised reading — fonds Alexandre Grothendieck, shelfmark 52, pages 1-13.
%
% This is an INTERPRETATION, not a transcription. Everything the critical
% apparatus carried moves into footnotes. In any dispute of substance the
% transcription governs: whoever cites Grothendieck must cite batch-01.fr.tex.
%
% Pass: Opus 5 (claude-opus-5), 2026-09-13 — first pass, unchecked against the
% pages by a human. The folder was read whole; it holds one batch, pages 1-13.
% Page 1 is a cover sheet carrying only the title, and takes no \pagerange.
%
% What the folder is. The Torelli theorem for a family of curves, written as a
% statement about a moduli functor over a base S: given a principally
% polarised abelian scheme A/S, the functor F of curves C/T with an
% isomorphism of polarised jacobians J^0_{C/T} -> A_T. Three movements:
% injectivity of Aut(C) -> Aut(J) and its corollaries (pages 2-4); the
% representability of F through a rigidified F' and a quotient by A (pages
% 4-7); the infinitesimal Torelli theorem, reduced to Max Noether's theorem
% (pages 8-9); and Matsusaka's criterion, which controls what a family of
% jacobians can specialise to (pages 10-13).
%
% Notation fixed for the whole document. J^0_{C/S} = Pic^0_{C/S}, J^1_{C/S}
% the Pic^0-torsor of degree-1 classes, in which C sits by the Abel map. His
% doubly underlined sheaves (T, V, N, Omega) are written plainly. theta or
% Theta the polarisation divisor; n = g = dim A throughout (he writes n on
% pages 11-12 and g on page 12's Cor. 1). omega = Omega^1_{C/k}.
%
% Substantive departures from the pages, each footnoted where it occurs:
% — page 3: the Lefschetz number of an automorphism acting trivially on
%   H^1 is 1 - 2g + 1 = 2 - 2g, negative; the page writes -1 + 2g - 1 =
%   2(g-1) > 0. What the argument needs is that it is non-zero, which holds
%   either way; the reading gives the right sign.
% — page 3, the Lemme writes « induit par s » for the automorphism induced by
%   sigma; and page 3's Cor. 2 writes a section « sur A » before A exists.
% — page 9: H^0(omega) (x) H^0(omega) -> H^0(omega^2) is surjective iff C is
%   non-hyperelliptic OR g = 2; the page says « ssi non hyperelliptique »,
%   which is false in genus 2, where every curve is hyperelliptic and both
%   sides have dimension 3.
% — page 13, Cor. 3 (ii): an isomorphism u compatible with phi and phi'
%   exists in general only up to the sign of phi; it exists and is unique on
%   the nose when the fibres are hyperelliptic, which is the case the page
%   states. The reading says why (the hyperelliptic involution acts as -1).
% — page 7: the quotient M'/A is announced to exist « universellement »; the
%   reading records that a free action gives it as an algebraic space and
%   that the leaf does not say more.
%
% Modern names given and footnoted as ours: Torelli morphism of stacks
% M_g -> A_g and its degree 2 onto its image (Oort-Steenbrink for the local
% statement); Max Noether's theorem; Lefschetz fixed-point formula; the
% Matsusaka-Hoyt-Ran criterion (the page names Matsusaka only); the
% Kodaira-Spencer map.
%
% What the folder announces and never establishes: the representability of
% F' (« rédiger proprement »); the existence of the quotient; the argument of
% the bottom two-thirds of page 10 (cancelled and illegible); the proofs of
% the corollaries of pages 12-13.
\documentclass[11pt,a4paper]{article}
% Le préambule commun à toutes les transcriptions du fonds Grothendieck.
%
% Il tient en une idée : **l'incertitude se compose, elle ne se résout pas.**
% Une transcription qui lisse un mot douteux a détruit la seule chose pour
% laquelle on la faisait — un lecteur doit pouvoir savoir, à chaque endroit,
% ce qui était sur le feuillet et ce qui a été ajouté. Les six macros qui
% suivent sont donc l'appareil critique, et non de la décoration.
%
% Les mêmes commandes sont reconnues par `scripts/render.mjs`, qui produit la
% vue de lecture du site. Une macro ajoutée ici doit l'être aussi là-bas,
% faute de quoi le rendu échouera bruyamment — ce qui est voulu : un rendu
% silencieusement faux est pire qu'un rendu absent.

\ProvidesPackage{grothendieck}[2026/08/08 Transcriptions du fonds Grothendieck]

\usepackage{amsmath,amssymb,amsthm}
\usepackage{mathrsfs}
\usepackage{stmaryrd}
% Les diagrammes commutatifs sont partout dans ce fonds : c'est la langue
% même de la géométrie algébrique de Grothendieck, et une transcription qui
% les rendrait en prose ne transcrirait rien.
\usepackage{tikz-cd}
% Français par défaut — c'est la langue des feuillets — mais l'anglais est
% chargé aussi : la traduction et le résumé sont des documents anglais, et la
% typographie française (espaces avant les deux-points) y serait une faute.
% Ces documents basculent par \selectlanguage{english} en tête de corps.
\usepackage[english,french]{babel}
\usepackage{fontspec}
\usepackage{geometry}
\geometry{a4paper,margin=25mm,marginparwidth=28mm}
\usepackage{marginnote}
\usepackage{xcolor}
% ulem, not soul. `soul`'s \st and \ul are built on a character-by-character
% scan that XeLaTeX's font machinery breaks: the document fails with an
% undefined control sequence rather than degrading. ulem does the same two jobs
% and is Unicode-engine safe. `normalem` stops it hijacking \emph, which the
% transcriptions use for Grothendieck's own emphasis.
\usepackage[normalem]{ulem}
\usepackage{hyperref}
\hypersetup{colorlinks=true,linkcolor=black,urlcolor=black,pdfborder={0 0 0}}

% --- Métadonnées du lot -----------------------------------------------------
% Reprises telles quelles par le script de rendu, qui les affiche en tête de
% la vue de lecture. Elles fixent ce que le fichier prétend couvrir : sans
% elles, une transcription flotte sans point d'attache dans un fonds de seize
% mille feuillets.

\newcommand{\folder}[1]{\gdef\grothFolder{#1}}
\newcommand{\batch}[1]{\gdef\grothBatch{#1}}
\newcommand{\leaves}[2]{\gdef\grothFirst{#1}\gdef\grothLast{#2}}
\newcommand{\foldertitle}[1]{\gdef\grothFoldertitle{#1}}
\gdef\grothFolder{?}\gdef\grothBatch{?}\gdef\grothFirst{?}\gdef\grothLast{?}\gdef\grothFoldertitle{}

% --- Le résumé --------------------------------------------------------------
%
% Il ouvre la lecture modernisée, sous le titre « Résumé », et il est composé
% en petit. Ce n'est pas un ornement : c'est le seul passage du document qui
% ne soit pas des mathématiques, et un lecteur doit pouvoir le reconnaître
% avant de l'avoir lu — puis le sauter s'il connaît déjà le sujet, ou s'y
% arrêter s'il ne le connaît pas. Le filet qui le suit marque la reprise du
% corps.

\newenvironment{resume}
  {\small\setlength{\parskip}{0.45em}}
  {\par\medskip\hrule\medskip}

% --- Mots-clés ---------------------------------------------------------------
%
% En anglais, en clôture du résumé de la lecture modernisée : le vocabulaire
% moderne sous lequel chercher ce que les feuillets construisent. Le manifeste
% du site (`npm run manifest`) les extrait pour étiqueter la cote entière —
% cette ligne est la source unique des tags, il n'y a pas de fichier à côté.
\newcommand{\keywords}[1]{\par\smallskip\noindent
  {\footnotesize\textsc{Keywords} --- \textit{#1}}\par}

% --- L'appareil critique ----------------------------------------------------

% Le numéro de feuillet, en marge. C'est l'ancre qui permet de régler un
% désaccord sur une formule en regardant *un* feuillet plutôt qu'une chemise
% de six cents. Le site s'en sert aussi pour tourner les pages du fac-similé
% quand on fait défiler la transcription.
\newcommand{\leaf}[1]{%
  \marginnote{\footnotesize\color{gray!70}\textsf{#1}}%
  \label{leaf:#1}\ignorespaces}

% Illisible : on ne devine pas. Un mot inventé qui se lit comme les autres est
% le pire résultat possible.
\newcommand{\ill}{\textcolor{red!60!black}{[\,\ldots\,]}}

% Lecture douteuse : le mot est donné, et signalé comme incertain.
\newcommand{\uncertain}[1]{\uline{#1}}

% Ajout de l'éditeur — une lettre manquante, un mot sous-entendu.
\newcommand{\add}[1]{\textcolor{blue!50!black}{[#1]}}

% Biffé par Grothendieck lui-même. Ce qu'il a rayé fait partie du document :
% une rature dit souvent où la pensée a changé de direction.
\newcommand{\struck}[1]{\sout{#1}}

% Note du transcripteur, sur l'état matériel du feuillet.
\newcommand{\note}[1]{\par\smallskip{\small\color{gray!60!black}\textit{[#1]}}\par\smallskip}

% Note marginale de Grothendieck, distincte du corps du texte.
\newcommand{\marginal}[1]{\par\smallskip\noindent\hspace{1em}%
  {\small\color{gray!60!black}#1}\par\smallskip}

% --- Titre ------------------------------------------------------------------

\AtBeginDocument{%
  \begin{center}
    {\large\bfseries Fonds Alexandre Grothendieck --- cote n\textsuperscript{o}~\grothFolder}\\[2pt]
    {\itshape\grothFoldertitle}\\[4pt]
    {\small feuillets \grothFirst--\grothLast\ (lot \grothBatch)}\\[2pt]
    {\footnotesize\color{gray!60!black}Transcription établie d'après le fac-similé numérisé
     par l'Université de Montpellier.\\
     Les lectures incertaines sont soulignées, les illisibles notées [\,\ldots\,],
     les ajouts entre crochets.}
  \end{center}
  \vspace{1em}}

\endinput


\folder{52}
\pages{1}{13}
\dating{s.d.}
\watermark{Édition de démonstration\\interprétation personnelle de l'œuvre}
\foldertitle{Théorème de Torelli : notes manuscrites (s.d.) — lecture modernisée du dossier entier}

\begin{document}

\section*{Résumé}

\begin{resume}
À une courbe algébrique --- une surface de Riemann compacte, si l'on préfère ---
on associe un objet plus gros et plus régulier, sa \emph{jacobienne} : un
tore complexe de dimension égale au genre $g$ de la courbe, qui est en même
temps un groupe, et qui porte une structure supplémentaire, une
\emph{polarisation}, venue de la manière dont les lacets de la courbe
s'intersectent. Le théorème de Torelli dit que rien ne se perd dans ce
passage : la jacobienne polarisée détermine la courbe. C'est un peu comme
reconnaître un instrument à la seule forme de son spectre.

Ces feuillets ne redémontrent pas le théorème classique ; ils se demandent ce
qu'il devient \emph{en famille}. On se donne une famille d'objets du type
jacobienne --- un schéma abélien polarisé $A$ au-dessus d'une base $S$ --- et
l'on cherche l'espace $M$ qui paramètre les courbes dont la jacobienne est la
fibre de $A$. Le théorème de Torelli dit que les fibres de $M$ au-dessus de $S$
ont au plus un point, ou deux ; le dossier veut savoir comment ces points
varient.

Trois outils, un par mouvement. D'abord un argument de points fixes : un
automorphisme de la courbe qui ne bouge pas la jacobienne doit avoir un point
fixe, par la formule de Lefschetz, et il est alors l'identité. C'est ce qui
rend le problème rigide. Ensuite un calcul au premier ordre : si l'on déforme
infiniment peu la courbe sans déformer sa jacobienne, la déformation est-elle
triviale ? La réponse tient à un théorème du dix-neuvième siècle, celui de Max
Noether sur les formes différentielles, et elle est oui, sauf pour les
courbes hyperelliptiques --- celles qui sont des revêtements doubles de la
droite --- où quelque chose se replie. Enfin, pour savoir ce qui arrive à la
limite, quand une famille de jacobiennes dégénère, un critère de Matsusaka :
une variété abélienne principalement polarisée qui contient une courbe de la
bonne classe est une jacobienne, ou un produit de jacobiennes.

Le brouillon est rapide et souvent illisible, surtout aux pages 5 à 7 et 10 à
12. Ce qui s'y lit nettement est le squelette d'un traitement moderne du
morphisme de Torelli entre l'espace des modules des courbes et celui des
variétés abéliennes : injectif sur les automorphismes, immersif hors du lieu
hyperelliptique, et d'image fermée à des produits près.

\keywords{Torelli theorem, Jacobian, principally polarized abelian scheme, Lefschetz fixed-point formula, infinitesimal Torelli, Max Noether theorem, Matsusaka criterion, Hilbert scheme}
\end{resume}

\section*{Le fil du dossier, et les conventions}

\begin{enumerate}
\item[1.] \emph{Rigidité} (pages 2 à 4). $\mathrm{Aut}_{S}(C) \to
\mathrm{Aut}_{S\text{-gr}}(J^{0}_{C/S})$ est une immersion fermée ; une
translation de la jacobienne qui conserve la courbe est nulle ; les
isomorphismes de courbes s'injectent dans ceux des jacobiennes.
\item[2.] \emph{Le problème de modules} (pages 4 à 7). Le foncteur $F$ des
courbes munies d'un isomorphisme de leur jacobienne polarisée sur $A$ ; une
variante rigidifiée $F'$ représentable par un $M'$ quasi-projectif sur lequel
$A$ opère ; $M = M'/A$.
\item[3.] \emph{Torelli infinitésimal} (pages 8 et 9), réduit par dualité au
théorème de Max Noether.
\item[4.] \emph{Spécialisation et critère de Matsusaka} (pages 10 à 13), et les
corollaires qui en sortent sur l'image de $M \to S$ et sur les isomorphismes
de courbes au-dessus d'une base réduite.
\end{enumerate}

\emph{Conventions.} Une courbe sur $S$ est un $S$-schéma lisse et propre à
fibres géométriquement connexes de dimension 1, de genre $g \geq 2$.
$J^{0}_{C/S} = \mathrm{Pic}^{0}_{C/S}$ est sa jacobienne, munie de sa
polarisation principale canonique, et $J^{1}_{C/S}$ le torseur sous
$J^{0}_{C/S}$ des classes de degré 1, dans lequel $C$ se plonge par
l'application d'Abel.\footnote{Il écrit ces deux objets d'un même $J$ orné,
distingués par l'exposant, et double le soulignement des foncteurs et des
faisceaux ($\mathrm{Pic}$, $\mathrm{Aut}$, $\mathrm{Isom}$, $\mathrm{Hilb}$,
$T$, $V$, $N$, $\Omega$) ; on les écrit sans marque. Il appelle
$J^{0}_{C/S}$ le « schéma jacobien », suivi d'un adjectif lu « homogène » sans
certitude.} On note $\omega = \Omega^{1}_{C/k}$ sur un corps, $\theta$ ou
$\Theta$ le diviseur de polarisation, et $n = g$ la dimension de $A$ ; il écrit
$n$ aux pages 11 et 12 et $g$ au corollaire de la page 12.

\emph{Deux rigidifications.} La page 5 rigidifie par des différentielles
$h$-uples, la page 7 par le plongement de la courbe dans $A$, c'est-à-dire par
un point du schéma de Hilbert de $A$. Les feuillets ne disent pas comment
elles s'articulent, et on les garde distinctes.

\emph{Ce que le dossier annonce et n'établit pas} : la représentabilité de
$F'$, qu'une note marginale renvoie à plus tard (« rédiger proprement ») ;
l'existence du quotient $M'/A$ ; l'argument de la page 10, dont les deux tiers
sont annulés ; les démonstrations des corollaires des pages 12 et 13.

\pagerange{2}{2}
\section*{Automorphismes de la courbe et de sa jacobienne (page 2)}

\textbf{Théorème 1.} Soit $S$ localement noethérien et $C$ une courbe de genre
$g \geq 2$ sur $S$. Alors le morphisme de foncteurs
\[
  \mathrm{Aut}_{S}(C) \longrightarrow \mathrm{Aut}_{S\text{-gr}}(J^{0}_{C/S})
\]
est injectif.

Principe. Soit $\sigma$ un automorphisme de $C$ qui induit l'identité sur
$J^{0}_{C/S}$. L'automorphisme $\bar{\sigma}$ qu'il induit sur le torseur
$J^{1}_{C/S}$ commute alors à l'action de $J^{0}$, donc est la translation
$T_{s}$ par une section $s$ de $J^{0}_{C/S}$, et il faut voir que $s = 0$. Le
cas d'un corps algébriquement clos est le Lemme de la page 3 ; l'énoncé
relatif s'en déduit parce que les deux foncteurs sont non ramifiés sur $S$ ---
« pas d'automorphismes infinitésimaux ! ».

\textbf{Corollaire 1.} Le morphisme ci-dessus est un monomorphisme, et même une
immersion fermée.

En effet, pour $g \geq 2$, $\mathrm{Aut}_{S}(C)$ est fini et non ramifié sur
$S$, puisque $H^{0}(C_{s}, T_{C_{s}}) = 0$ en chaque point. Le noyau $K$ du
morphisme est un sous-groupe fermé de $\mathrm{Aut}_{S}(C)$, fini et non
ramifié sur $S$, dont les fibres géométriques sont réduites à l'élément neutre
par le Lemme ; un tel $K$ est une immersion fermée dans $S$ qui admet une
section, donc $K = S$. Un monomorphisme propre étant une immersion fermée,
le corollaire suit.\footnote{La page énonce le corollaire et en indique les
deux ingrédients --- $\mathrm{Aut}_{S}(C)$ « projectif sur $S$ » et les deux
schémas « non ramifiés » ---, mais la ligne qui les assemble est illisible ;
l'assemblage est le nôtre. Le passage « monomorphisme propre $\Rightarrow$
immersion fermée » est EGA~IV, 18.12.6. Deux lignes insérées puis biffées ne se
lisent pas.}

\pagerange{3}{3}
\section*{Le lemme des points fixes (page 3)}

\textbf{Lemme.} Soit $C$ une courbe projective lisse et connexe de genre
$g \geq 2$ sur un corps algébriquement clos $k$, et $\sigma$ un automorphisme
de $C$ qui induit l'identité sur $J^{0}_{C/k}$. Alors $\sigma$ est l'identité.

\emph{Démonstration.} L'automorphisme de $J^{1}$ induit par $\sigma$ est une
translation $T_{a}$, $a \in J^{0}(k)$, et pour tout $x \in C(k)$, vu dans
$J^{1}$, on a $\sigma(x) = x + a$.\footnote{La page écrit « l'automorphisme de
$J^{1}_{C/S}$ induit par $s$ » ; c'est $\sigma$ qui induit, $s$ ou $a$ étant
ce qu'on cherche. Elle garde aussi l'indice $C/S$ dans un énoncé posé sur un
corps.} Il suffit donc que $\sigma$ ait un point fixe $x$ : alors
$a = \sigma(x) - x = 0$. Or, pour $\ell$ premier à la caractéristique, $\sigma$
opère trivialement sur $H^{1}(C, \mathbb{Q}_{\ell})$, dual du module de Tate de
$J^{0}$, de dimension $2g$. Si $\sigma \neq \mathrm{id}$, le graphe de $\sigma$
et la diagonale de $C \times C$ sont deux courbes distinctes, et leur nombre
d'intersection est donné par la formule des traces de Lefschetz :
\[
  (\Gamma_{\sigma} \cdot \Delta)
  = \sum_{i=0}^{2} (-1)^{i}\, \mathrm{tr}\bigl(\sigma^{*} \mid H^{i}\bigr)
  = 1 - 2g + 1 = 2 - 2g \neq 0.
\]
Les deux courbes se rencontrent donc, $\sigma$ a un point fixe, et
$a = 0$.\footnote{La page écrit « $-1 + 2g - 1 = 2(g-1) > 0$ », et conclut à
l'existence d'un point fixe. Le signe est inversé : la trace sur $H^{1}$ entre
avec le signe moins, et le nombre de Lefschetz vaut $2 - 2g$, qui est
\emph{négatif}. Ce qui sert à l'argument est seulement qu'il n'est pas nul, et
cela reste vrai. On peut d'ailleurs conclure plus vite : pour
$\sigma \neq \mathrm{id}$ les points fixes sont isolés et comptent chacun avec
une multiplicité positive, de sorte qu'un nombre de Lefschetz négatif est
contradictoire par lui-même. Une première version du calcul, biffée sous
celle qui reste, ne se lit pas.}

\textbf{Corollaire 2.} Sous les hypothèses du Théorème 1, une section $s$ de
$J^{0}_{C/S}$ telle que la translation $T_{s}$ de $J^{1}_{C/S}$ conserve $C$
est la section nulle.

Sur les fibres géométriques, $T_{a}$ restreinte à $C$ est un automorphisme
dont le prolongement à $J^{1}$ est une translation, donc qui opère
trivialement sur $J^{0}$ ; par le Lemme il est l'identité, et $a = 0$. On
conclut sur $S$ comme au Corollaire 1.\footnote{La page écrit « une section
$s$ de $J^{0}_{C/S}$ sur $A$ », avec un $A$ majuscule, alors que $A$ ne sera
introduit qu'à la page 4 ; on lit $S$. « Section nulle » est écrit en
interligne au-dessus de « l'identité » biffé.}

C'est ce corollaire qui servira à la page 7 : il dit que la translation par
$A$ opère \emph{librement} sur les courbes tracées dans $A$.

\pagerange{4}{4}
\section*{Isomorphismes, et le problème de modules (page 4)}

Le haut de la page reprend le Corollaire 2 au-dessus des nombres duaux
$k[\varepsilon]$, puis d'une extension de carré nul
$\mathcal{O}_{S} \oplus \mathcal{I}$, pour en tirer la version infinitésimale ;
le bloc est encadré puis annulé.\footnote{Deux longues diagonales ; une
partie seulement se lit, et un « Prop. » biffé le suit.}

\textbf{Corollaire 3.} Soient $C$, $C'$ deux courbes de genre $g \geq 2$ sur
$S$. Alors
\[
  \mathrm{Isom}_{S}(C, C') \longrightarrow
  \mathrm{Isom}_{S\text{-gr}}(J^{0}_{C/S}, J^{0}_{C'/S})
\]
est une immersion fermée. C'est le Corollaire 1 transporté aux torseurs : là où
$\mathrm{Isom}_{S}(C, C')$ n'est pas vide, il est localement un torseur sous
$\mathrm{Aut}_{S}(C)$.

\emph{Le problème de modules.} Soit $A$ un schéma abélien de dimension relative
$g$ sur $S$, \emph{polarisé au sens strict} : on se donne une section du schéma
de Néron-Severi.\footnote{« Nér.-Sév. » sur la page. Pour que la question ait
un sens, la polarisation doit être principale, puisque celle d'une jacobienne
l'est ; c'est ce que dit la condition $\deg(\Theta^{n})/n! = 1$ des pages 11 et
12, et on la sous-entend dès ici.} Pour tout $S$-schéma $T$, soit $F(T)$
l'ensemble des classes d'isomorphisme de couples $(C, \varphi)$, où $C$ est une
courbe de genre $g$ sur $T$ et
\[
  \varphi : J^{0}_{C/T} \xrightarrow{\;\sim\;} A_{T}
\]
un isomorphisme compatible avec les polarisations.

\emph{Ce que dit le théorème classique.} Sur un corps algébriquement clos, si
$A = J^{0}_{C_{0}}$, le théorème de Torelli dit que toute $C$ de $F(k)$ est
isomorphe à $C_{0}$, et $F(k)$ s'identifie à
$\mathrm{Aut}(A, \theta)/\mathrm{Aut}(C_{0})$. Or
$\mathrm{Aut}(J^{0}_{C_{0}}, \theta) = \mathrm{Aut}(C_{0}) \times \{\pm 1\}$
si $C_{0}$ n'est pas hyperelliptique, tandis que $-1$ provient de l'involution
hyperelliptique si elle l'est. Ainsi $F(k)$ a \emph{deux} points au-dessus
d'une jacobienne non hyperelliptique, et un seul au-dessus d'une jacobienne
hyperelliptique.\footnote{Ce calcul est le nôtre ; il n'est pas sur la page 4.
Il éclaire deux fragments illisibles de la suite : le « $\mathbb{Z}/2$ » de la
page 6, et, page 9, la fibre de $M$ réduite à $\mathrm{Spec}\,k(s)$ dans le cas
hyperelliptique. Dans la langue d'aujourd'hui, $F$ est la fibre du morphisme
de Torelli $\mathcal{M}_{g} \to \mathcal{A}_{g}$ entre champs de modules
au-dessus du point de $\mathcal{A}_{g}$ défini par $A$, et ce morphisme est de
degré 2 sur son image, ramifié le long du lieu hyperelliptique.}

\pagerange{5}{7}
\section*{Rigidification, schéma de Hilbert, quotient (pages 5 à 7)}

Ces trois pages sont écrites très vite et largement biffées ; on n'en donne
que les articulations qui se lisent.

\emph{Une première rigidification} (page 5). $F'(T)$ est l'ensemble analogue à
$F(T)$ où l'on ajoute une rigidification par des éléments de
$H^{0}(C, \omega^{\otimes h})$, $h \geq 3$. Le nombre qui intervient est
\[
  v = 2h(g-1) - (g-1) = (2h-1)(g-1) = \dim H^{0}(C, \omega^{\otimes h}),
\]
par le théorème de Riemann-Roch, puisque $\deg \omega^{\otimes h} = 2h(g-1)$
dépasse $2g - 2$.\footnote{L'interprétation de $v$ comme dimension est la
nôtre ; la page pose $v$ sans dire ce qu'il compte. Pour $h \geq 3$, le
faisceau $\omega^{\otimes h}$ est très ample, ce qui explique la borne. Un bloc
en biais dans la marge gauche mentionne des courbes hyperelliptiques, sans
plus.}

\textbf{Théorème (annoncé).} Le foncteur $F'$ est représentable par un $M'$
quasi-projectif sur $S$.\footnote{Suivi de fragments sur la relation
d'équivalence $F' \times_{F} F' \rightrightarrows F'$ et de la mention
« rédiger proprement » en marge. Rien de la démonstration ne se lit.}

\emph{Le cas d'un corps} (page 6). Pour $S = \mathrm{Spec}(k)$, la page renvoie
à la théorie de Hilbert et au théorème de Torelli « sous sa forme classique »,
et parle d'un « $\mathbb{Z}/2$ » ; le dénombrement de la section précédente dit
ce qu'il peut être.\footnote{La page est de la même farine que la précédente :
les mots « Hilbert », « Torelli », « classique », « théorème » portent, la
prose presque pas. Le $\mathbb{Z}/2$ est écrit « $\mathbb{Z}/2 \in A_{-g}(A)$ »,
formule que cette lecture ne sait pas interpréter.}

\emph{Le quotient} (page 7). Le schéma abélien $A$ opère sur $F'$ par
translation, et, d'après le Théorème 1 --- en fait son Corollaire 2 ---, il
opère librement. Si $F'$ est représentable par $M'$, et si le quotient $M'/A$
existe « universellement », alors $M = M'/A$ représente $F$ ; $A$ est propre
et plat sur $S$.\footnote{Une action libre d'un schéma en groupes propre et
plat donne un quotient au moins comme espace algébrique ; qu'il soit un schéma
demande davantage, et la page, qui met « universellement » entre guillemets,
ne dit pas quoi.}

\emph{La seconde rigidification.} $F'(T)$ s'explicite comme une partie de
$\mathrm{Hilb}_{A/S}(T)$ : une courbe $C$ tracée dans $A_{T}$ par un plongement
qui induit l'isomorphisme naturel $J^{0}_{C/T} \to A_{T}$. La page invoque le
théorème de Hilbert --- la projectivité du schéma de Hilbert de $A$, polarisé
--- et le nom de Matsusaka, pour savoir quels sous-schémas de $A$ sont de telles
courbes.\footnote{Le passage est une suite de fragments autour de
« Hilb », « plongement $C \to A_{T}$ », « Matsusaka », « théorème de Hilbert »,
« $\Theta$ » ; l'enchaînement est le nôtre, et il est celui que les pages 11 et
12 suivent.}

\pagerange{8}{9}
\section*{Le théorème de Torelli infinitésimal (pages 8 et 9)}

Soit $C$ une courbe sur un corps algébriquement clos $k$, plongée dans sa
jacobienne $A$. La question est : \emph{toute déformation du premier ordre de
$C$ dans $A$ provient-elle d'une translation infinitésimale de $A$ ?} Une telle
translation est d'ailleurs bien déterminée, puisque $C$ n'a pas
d'automorphismes infinitésimaux.

Les déformations de $C$ dans $A$ sont les éléments de $H^{0}(C, N)$, où $N$
est le faisceau normal. On a la suite exacte
\[
  0 \longrightarrow T \longrightarrow V \longrightarrow N \longrightarrow 0,
\]
où $T$ est le faisceau tangent de $C$, de degré $2 - 2g < 0$, et
$V = T_{A}|_{C} \simeq \mathcal{O}_{C}^{g}$ la restriction du faisceau tangent
de $A$, qui est libre. D'où
\[
  H^{0}(C, T) = 0, \qquad
  H^{0}(C, V) \simeq k^{g} \simeq H^{0}(A, T_{A}),
\]
et la suite longue
\[
  0 \to H^{0}(C, V) \to H^{0}(C, N) \to H^{1}(C, T)
    \to H^{1}(C, V) \simeq H^{1}(A, T_{A}).
\]
Les translations sont l'image de $H^{0}(C, V)$ ; la question est donc la
surjectivité de $H^{0}(C, V) \to H^{0}(C, N)$, c'est-à-dire l'injectivité de
\[
  H^{1}(C, T) \longrightarrow H^{1}(A, T_{A}).
\]
La source classe les déformations du premier ordre de $C$, le but celles de
$A$, et la flèche envoie une déformation de la courbe sur celle qu'elle induit
de sa jacobienne.\footnote{C'est l'application de Kodaira-Spencer du morphisme
de Torelli ; le nom n'est pas sur la page. L'isomorphisme
$H^{1}(C, V) \simeq H^{1}(A, T_{A})$ vient de
$H^{1}(C, \mathcal{O}_{C}) \simeq H^{1}(A, \mathcal{O}_{A})$.} D'où l'énoncé :

\textbf{Corollaire.} Si $C$ n'est pas hyperelliptique, ou si $g = 2$, toute
déformation du premier ordre de $C$ qui induit une déformation triviale de
$J^{0}_{C}$ est triviale.

\emph{Démonstration.} Par la dualité de Serre, $H^{1}(C, T)^{\vee} =
H^{0}(C, \omega^{\otimes 2})$ et $H^{1}(C, V)^{\vee} = H^{0}(C, V^{\vee}
\otimes \omega) = H^{0}(C, \omega) \otimes H^{0}(C, \omega)$. L'injectivité
cherchée équivaut donc à la surjectivité de la multiplication
\[
  H^{0}(C, \omega) \otimes H^{0}(C, \omega) \longrightarrow
  H^{0}(C, \omega^{\otimes 2}).
\]
Pour $g \geq 3$, c'est le théorème de Max Noether : elle est surjective si et
seulement si $C$ n'est pas hyperelliptique. Pour $g = 2$ elle est toujours
surjective.\footnote{La page écrit « vrai \emph{ssi} $C$ est non
hyperelliptique », ce qui est faux en genre 2 : toute courbe y est
hyperelliptique, et pourtant l'image de la multiplication, qui est
$\mathrm{Sym}^{2} H^{0}(\omega)$, de dimension 3, remplit
$H^{0}(\omega^{\otimes 2})$, de dimension $3g - 3 = 3$. Pour $C$
hyperelliptique de genre $g \geq 3$, l'image est de dimension $2g - 1$ et le
conoyau de dimension $g - 2$. Le nom de Max Noether n'est pas sur la page.}

\emph{Le cas hyperelliptique.} Le noyau de $H^{1}(C, T) \to H^{1}(A, T_{A})$
n'est alors pas nul, et la page en tire que la fibre de $M$ au-dessus du point
correspondant se réduit à $\mathrm{Spec}\,k(s)$ : les deux points que $M$ a
au-dessus d'une jacobienne générale, $(C, \varphi)$ et $(C, -\varphi)$, s'y
confondent, et $M \to S$ y est ramifié.\footnote{La phrase est lacunaire ; on
lit « de dimension 1 » et « la fibre de $M$ en $s$ … $\mathrm{Spec}\,k(s)$ »
entre des mots illisibles. Si la dimension 1 porte sur le noyau ci-dessus,
elle ne vaut que pour $g = 3$. Le résultat moderne correspondant --- le
morphisme de Torelli entre espaces de modules grossiers est une immersion hors
du lieu hyperelliptique, et le morphisme de champs n'y est pas immersif --- est
dû à Oort et Steenbrink (1980) ; on ne sait pas si les feuillets sont
antérieurs, faute de datation.}

\pagerange{9}{10}
\section*{Relever une courbe (pages 9 et 10)}

La fin de la page 9 et le haut de la page 10 tirent de ce qui précède un énoncé
de relèvement : soit $s \in S$, et $C_{0}$ une courbe sur $k(s)$ munie d'un
isomorphisme compatible aux polarisations
\[
  J^{0}_{C_{0}/k(s)} \xrightarrow{\;\sim\;} A_{s}.
\]
Il s'agit de prolonger $C_{0}$ en une courbe $C$ sur $S$ munie d'un
isomorphisme $J^{0}_{C/S} \simeq A$, et de façon unique. C'est ce que donne un
morphisme $M \to S$ non ramifié au point considéré, par le Corollaire des pages
8 et 9.\footnote{Les deux énoncés (ii) et « Corollaire » de la page 9 sont
presque entièrement illisibles, de même que le haut de la page 10 ; les deux
tiers inférieurs de celle-ci sont un bloc encadré et annulé par une nappe de
diagonales, qui n'est pas lisible. Ce qui est donné est ce que les formules
lisibles et le mot « unique » permettent ; le lien avec la non-ramification est
le nôtre. Une note en biais dans la marge gauche de la page 9 ne livre que le
mot « finalement ».}

\pagerange{11}{12}
\section*{Le critère de Matsusaka (pages 11 et 12)}

Il faut maintenant savoir ce qui arrive à la limite : quand une famille de
courbes $C_{y}$ tracées dans $A$ se spécialise en un cycle $C_{a}$, celui-ci
est-il encore une courbe, et $A_{a}$ encore une jacobienne ?

\textbf{Critère de Matsusaka.} Soit $A$ une variété abélienne de dimension $n$
sur un corps algébriquement clos, $\Theta$ un diviseur positif définissant une
polarisation principale, c'est-à-dire $\deg(\Theta^{n})/n! = 1$, et $Z$ un
1-cycle positif tel que
\[
  \Theta^{\,n-1} \equiv (n-1)!\, Z .
\]
Alors $Z$ est de multiplicité 1, et $(A, \Theta)$ est le produit, comme
variété abélienne principalement polarisée, des jacobiennes des composantes de
$Z$ ; en particulier, si $Z$ est irréductible, $(A, \Theta)$ est la jacobienne
de la courbe $Z$.\footnote{La page attribue le résultat à Matsusaka, nom lu
avec un doute léger, et en énonce la forme irréductible ; les mots
« multiplicité 1 » et « produit des jacobiennes des composantes de $C_{a}$ »,
dans le fragment qui précède, montrent qu'il en avait aussi la forme
réductible. Historiquement : T.~Matsusaka, 1959, pour le cas d'une courbe
irréductible ; W.~Hoyt, 1963, pour les produits ; Z.~Ran, 1981, pour une
forme plus générale. On ne sait pas lesquels les feuillets connaissent.}

\emph{Application.} Dans la famille, les classes se conservent :
\[
  \frac{\deg \Theta_{y}^{n}}{n!} = \frac{\deg \Theta_{a}^{n}}{n!} = 1,
  \qquad
  \Theta_{y}^{\,n-1} \equiv (n-1)!\, C_{y},
  \qquad
  \Theta_{a}^{\,n-1} \equiv (n-1)!\, C_{a}.
\]
Le cycle limite $C_{a}$ satisfait donc aux hypothèses du critère : il est
réduit, $(A_{a}, \Theta_{a})$ est un produit de jacobiennes, et c'est une
jacobienne dès que $C_{a}$ est irréductible, ce qui arrive exactement quand
$\Theta_{a}$ est irréductible.\footnote{La formule est écrite deux fois à la
page 11, la première biffée ; la classe que l'énoncé général appelle $Z$ est ici
la courbe elle-même. La conservation des classes dans une famille plate de
cycles, et l'existence de la limite par propreté du schéma de Hilbert, sont
sous-entendues. L'équivalence « $C_{a}$ irréductible si et seulement si
$\Theta_{a}$ irréductible » est ce que le haut de la page 12, barré, paraît
dire ; on la tire du critère lui-même : un produit non trivial de variétés
abéliennes principalement polarisées a un diviseur thêta réductible, et une
jacobienne l'a irréductible.}

\textbf{Corollaire 1.} Soit $f : M \to S$ le morphisme structural, et
$T = \overline{f(M)}$. Un point $t \in T$ est dans $f(M)$ si et seulement si le
diviseur de la polarisation de $A_{t}$ --- défini à translation près,
$\deg \Theta^{g}/g! = 1$ --- est irréductible.

Autrement dit : l'adhérence du lieu des jacobiennes est faite des produits de
jacobiennes, et l'on reconnaît une jacobienne parmi eux à l'irréductibilité de
son diviseur thêta.

\pagerange{13}{13}
\section*{Isomorphismes au-dessus d'une base réduite (page 13)}

Un premier corollaire, qui comparait deux courbes $C$, $C'$ munies
d'isomorphismes $\varphi : J^{0}_{C/S} \to A$ et
$\varphi' : J^{0}_{C'/S} \to A$, est encadré puis annulé ; il est repris sous
le numéro 3.

\textbf{Corollaire 3.} Soit $S$ réduit.

\begin{enumerate}
\item[(i)] La page énonce une condition nécessaire et suffisante pour que $A$
soit une jacobienne polarisée, portant sur ses fibres et sur l'irréductibilité
des diviseurs de Matsusaka.\footnote{L'énoncé (i) est très lacunaire ; il porte
une hypothèse sur les points où $A_{s}$ est « hyperelliptique » et une
condition sur les composantes de $S$, dont on ne lit pas l'articulation. On
n'en donne que la forme, et cette lecture ne le garantit pas.}
\item[(ii)] Supposons qu'en tout point $s \in S$, $A_{s}$ soit la jacobienne
d'une courbe hyperelliptique. Soient $C$, $C'$ deux courbes sur $S$ munies
d'isomorphismes de jacobiennes polarisées
\[
  \varphi : J^{0}_{C/S} \longrightarrow A,
  \qquad
  \varphi' : J^{0}_{C'/S} \longrightarrow A.
\]
Alors il existe un $S$-isomorphisme $u : C \to C'$ compatible avec $\varphi$ et
$\varphi'$, et un seul.
\end{enumerate}

L'hypothèse hyperelliptique est ce qui rend l'énoncé exact. Par le
Corollaire 3 de la page 4, $\mathrm{Isom}_{S}(C, C') \to
\mathrm{Isom}_{S\text{-gr}}(J^{0}_{C/S}, J^{0}_{C'/S})$ est une immersion
fermée ; son image réciproque par la section $\varphi'^{-1} \circ \varphi$ est
donc un sous-schéma fermé $Z$ de $S$, qui représente les $u$ compatibles, et
qui en a au plus un. Un point $s$ est dans $Z$ dès qu'il existe, sur une
clôture algébrique de $k(s)$, un isomorphisme $u_{s}$ avec
$\varphi'_{s} \circ J(u_{s}) = \varphi_{s}$. Le théorème de Torelli classique
donne un $u_{s}$ avec $\varphi'_{s} \circ J(u_{s}) = \pm \varphi_{s}$, et quand
$-1$ est induit par l'involution hyperelliptique $\iota$ on corrige le signe en
remplaçant $u_{s}$ par $u_{s} \circ \iota$. Ainsi $Z$ contient tous les points
de $S$, et $Z = S$ puisque $S$ est réduit.\footnote{Sans l'hypothèse
hyperelliptique, (ii) est faux : pour $C$ non hyperelliptique,
$(C, \varphi)$ et $(C, -\varphi)$ ne sont pas isomorphes, et il n'existe aucun
$u$ compatible avec $\varphi$ et $-\varphi$. L'hypothèse est à demi lisible
sur la page (« … que $A_{s}$ soit \emph{hyperelliptique} ») et l'on lui donne
cette forme. La réduction de $S$ est exactement ce qui fait passer d'un
sous-schéma fermé qui contient tous les points à $S$ lui-même. L'argument est
le nôtre ; le feuillet s'arrête sur l'énoncé.}

\end{document}
